% Acesso a um registrador do GPIO: rdy registrado, um ciclo após o vld.
\begin{tikzpicture}[font=\scriptsize]
  \def\w{1.1}
  \foreach \y/\n in {0/clk,-0.8/vld\_i,-1.6/{addr\_i, we\_i, data\_i},-2.4/rdy\_o,-3.2/{data\_o ou registrador}} \node[font=\tiny\ttfamily,anchor=east] at (-0.1,\y-0.2) {\n};
  \foreach \i in {0,...,5} \draw[thick] ({\i*\w},-0.4) -- ({\i*\w},0) -- ({\i*\w+\w/2},0) -- ({\i*\w+\w/2},-0.4) -- ({\i*\w+\w},-0.4);
  \draw[thick,azul] (0,-1.2) -- ({1*\w},-1.2) -- ({1*\w},-0.8) -- ({3*\w},-0.8) -- ({3*\w},-1.2) -- ({6*\w},-1.2);
  \draw[thick] (0,-1.8) -- ({1*\w-0.08},-1.8) -- ({1*\w},-1.6) -- ({3*\w},-1.6) -- ({3*\w+0.08},-1.8) -- ({6*\w},-1.8);
  \draw[thick] ({1*\w-0.08},-1.8) -- ({1*\w},-2.0) -- ({3*\w},-2.0) -- ({3*\w+0.08},-1.8);
  \node[font=\tiny] at ({2*\w},-1.8) {estáveis};
  \draw[thick,laranja!80!black] (0,-2.8) -- ({2*\w},-2.8) -- ({2*\w},-2.4) -- ({3*\w},-2.4) -- ({3*\w},-2.8) -- ({6*\w},-2.8);
  \draw[thick] (0,-3.4) -- ({2*\w-0.08},-3.4) -- ({2*\w},-3.2) -- ({6*\w},-3.2);
  \draw[thick] ({2*\w-0.08},-3.4) -- ({2*\w},-3.6) -- ({6*\w},-3.6);
  \node[font=\tiny] at ({4*\w},-3.4) {atualizado};
  \begin{pgfonlayer}{background}\fill[green!14] ({2*\w},-3.75) rectangle ({3*\w},0.15);\end{pgfonlayer}
  \foreach \x/\t in {1/1,2/2,3/3} \node[draw,circle,inner sep=0.8pt,font=\tiny\bfseries,fill=white] at ({\x*\w},0.45) {\t};
  \node[font=\tiny,text=corpo,align=left,anchor=north west] at ({6.2*\w},0.2) {%
    1. o mestre ergue \texttt{vld\_i} com endereço\\\quad e dados estáveis;\\
    2. na borda seguinte, o GPIO registra \texttt{rdy\_o}\\\quad e atualiza a saída ou o registrador;\\
    3. o mestre consome o dado e baixa \texttt{vld\_i};\\\quad \texttt{rdy\_o} volta a 0.};
\end{tikzpicture}
